%% ========================================================================
%% Chapter 2: Hardware Management and Basic Commands
%% ========================================================================

\chapter{Hardware Management and Basic Commands}
\label{ch:hardware-and-basic-commands}

\chapterhero{teal}{Learn how Linux detects hardware, loads drivers, and provides basic terminal commands so you can begin working with confidence.}{\faMicrochip\quad hardware}{\faCogs\quad driver}{\faTerminal\quad terminal}

Linux does not simply assume hardware exists. When the computer starts, the kernel checks each available bus, identifies attached devices, and loads the appropriate \term{driver}. The processor, RAM, network card, and disk all must be recognized before the system can be used. In Windows, this summary is available in Device Manager. In Linux, you can read the full trace directly from the terminal, often with much deeper detail.

Linux detects hardware through a staged process. When the computer starts, the kernel scans buses such as PCI and USB, then matches detected devices with suitable kernel modules. You can inspect the detection results with commands such as \cmd{lscpu}, \cmd{lspci}, \cmd{lsusb}, and \cmd{dmesg}. The process is automatic, but its traces remain auditable from the terminal.

Understanding hardware and basic terminal commands is the foundation for every later Linux lab. Without these skills, you will not be able to read system information, move between directories confidently, manage simple files, or interpret command documentation throughout this book.

\textbf{What You'll Learn}

By the end of this chapter, you will be able to:

\begin{enumerate}
    \item identify hardware components using \cmd{lspci}, \cmd{lsusb}, \cmd{lscpu}, and \cmd{dmesg};
    \item explain the function of kernel modules and view active modules with \cmd{lsmod} and \cmd{modinfo};
    \item use basic terminal commands for navigation and file and directory operations;
    \item read file contents and command output with \cmd{cat}, \cmd{head}, and \cmd{tail};
    \item filter and process text using \cmd{grep}, \cmd{sed}, and \cmd{awk};
    \item use built-in documentation such as \cmd{man}, \cmd{--help}, and \cmd{info} to learn new commands independently.
\end{enumerate}

\begin{notebox}
All labs in this chapter use the directory \dir{~/lab-os/chapter02-hardware/}.
Create this directory before starting the first lab:
\begin{lstlisting}[language=bash]
mkdir -p ~/lab-os/chapter02-hardware
cd ~/lab-os/chapter02-hardware
\end{lstlisting}
\end{notebox}

%% ------------------------------------------------------------------------
\section{Hardware Detection}
\label{sec:hardware-detection}

Whenever Linux starts, the kernel checks all attached devices. The results are stored and can be read at any time. The \cmd{lscpu} command displays processor architecture, physical core count, thread count, and virtualization features. The \cmd{lspci} command shows all devices on the PCI bus, including network cards, graphics cards, and storage controllers. For USB devices, use \cmd{lsusb}. If you need a more complete hardware report in one command, \cmd{lshw} is a good choice, although it requires superuser privileges.

Kernel messages produced during boot are stored in a ring buffer that can be read with \cmd{dmesg}. There you can see device detection records: when a disk was found, which \term{driver} was loaded, and whether a device failed to initialize. Add the \texttt{-T} flag to display time in a human-readable format.

\begin{tipbox}
Some hardware detection commands may not be installed in a minimal installation. Install the required packages with:
\begin{lstlisting}[language=bash]
sudo apt update
sudo apt install -y pciutils usbutils lshw
\end{lstlisting}
\end{tipbox}

\subsection*{Lab 2.1: Identify VM Hardware Components}

The goal of this lab is to read processor, memory, and PCI device specifications from a \term{virtual machine}.

\begin{enumerate}
    \item Display processor information:
\begin{lstlisting}[language=bash, caption={Reading processor specifications}]
lscpu
\end{lstlisting}
    \textit{Observe:} Notice the \texttt{CPU(s)}, \texttt{Core(s) per socket}, and \texttt{Vendor ID} lines. A virtual machine usually shows the host processor with a limited number of cores.

    \item View the list of PCI devices:
\begin{lstlisting}[language=bash, caption={PCI device list with active drivers}]
lspci -nnk
\end{lstlisting}
    \textit{Notice:} Each entry shows the vendor and device ID in brackets, followed by a \texttt{Kernel driver in use} line that shows the active \term{driver} module name.

    \item Find the network card and the \term{driver} that handles it:
\begin{lstlisting}[language=bash, caption={Filtering network card information}]
lspci -nnk | grep -A3 -i ethernet
\end{lstlisting}

    \item View kernel messages from the boot process:
\begin{lstlisting}[language=bash, caption={Kernel messages during boot}]
dmesg -T | head -50
\end{lstlisting}
    \textit{Observe:} Look for lines that mention device names such as \texttt{sda}, \texttt{eth0}, or \texttt{ens3}. Those lines mark when the kernel first detected the device.

    \item Search boot messages related to virtualization or storage devices:
\begin{lstlisting}[language=bash, caption={Filtering kernel messages with dmesg}]
dmesg -T | grep -Ei "virtio|nvme|sata|usb"
\end{lstlisting}
\end{enumerate}

\begin{challengebox}
Choose one active kernel module from the \cmd{lsmod} output. Use \cmd{modinfo} to find its description, license, and module filename. Then use \cmd{dmesg -T | grep -i module-name} to check whether any kernel messages mention that module during system startup.
\end{challengebox}

%% ------------------------------------------------------------------------
\section{Kernel Modules and Device Drivers}
\label{sec:kernel-modules}

The Linux kernel does not load all \term{driver} code at once. Instead, drivers are available as \textit{kernel modules}, which are code files that can be loaded when needed and removed when no longer required. The idea is like plugging in an electrical device only when it is needed, rather than leaving it connected all the time. Modules are stored in \dir{/lib/modules/\$(uname -r)/} and usually have the \texttt{.ko} extension.

The \cmd{lsmod} command shows all currently active modules, along with their size and usage count. To view details for one module, use \cmd{modinfo module-name}: you will see its description, available parameters, author name, license, and module file path. The \cmd{modprobe module-name} command loads a module and its dependencies automatically.

When the kernel successfully loads a module, the result is usually visible in \cmd{dmesg} messages, PCI/USB device listings, or new capabilities available on the system. For that reason, the three tools most often used together are \cmd{lsmod} to see what is active, \cmd{modinfo} to understand a specific module, and \cmd{dmesg} to trace when the \term{driver} was used or caused an error.

\begin{warningbox}
Do not remove a module that is being used by the active system. Removing a disk \term{driver} or network module can make the system unresponsive or disconnect an SSH session.
\end{warningbox}

\subsection*{Lab 2.2: Explore Active Modules and Device Drivers}

The goal of this lab is to view running kernel modules and connect them to devices detected by the system.

\begin{enumerate}
    \item Display the running kernel version:
\begin{lstlisting}[language=bash, caption={Active kernel version}]
uname -r
\end{lstlisting}

    \item View the list of active modules:
\begin{lstlisting}[language=bash, caption={Active kernel module list}]
lsmod | head -20
\end{lstlisting}
    \textit{Observe:} The first column is the module name, the second column is the size in bytes, and the third column is the number of processes using the module.

    \item Read details for the \texttt{loop} module, which creates virtual disks from files:
\begin{lstlisting}[language=bash, caption={Loop module details}]
modinfo loop
\end{lstlisting}

    \item Make sure the \texttt{loop} module is loaded, then verify it:
\begin{lstlisting}[language=bash, caption={Load and verify the loop module}]
sudo modprobe loop
lsmod | grep loop
\end{lstlisting}

    \item See how \cmd{modprobe} resolves module dependencies without actually reloading the module:
\begin{lstlisting}[language=bash, caption={Simulating module loading and dependencies}]
modprobe -n -v loop
\end{lstlisting}
    \textit{Notice:} The \texttt{-n} flag performs a simulation, while \texttt{-v} displays the steps that would be taken. This is useful when you want to understand module dependencies without changing the system state.

    \item Connect modules with real devices using PCI output and kernel messages:
\begin{lstlisting}[language=bash, caption={Connecting modules, devices, and kernel messages}]
lspci -nnk | grep -A3 -i -E "ethernet|network|vga|storage"
dmesg -T | grep -Ei "virtio|nvme|usb|ethernet" | tail -20
\end{lstlisting}
\end{enumerate}

\begin{challengebox}
Choose one device from the \cmd{lspci -nnk} output. Record the device name and its active \term{driver}. Then use \cmd{modinfo} to find the module description, and verify with \cmd{dmesg} whether any boot message mentions the same device or \term{driver}.
\end{challengebox}

%% ------------------------------------------------------------------------
\section{Basic Terminal Commands}
\label{sec:basic-commands}

The Linux terminal is the main interface for server administration. Unlike a graphical interface that hides details, the terminal gives full control: you know exactly what happens and why. Mastering basic commands is not optional. It is required.

Navigation starts with three commands: \cmd{pwd} displays the current directory location, \cmd{ls} shows directory contents, and \cmd{cd} moves to another directory. The \texttt{-l} flag for \cmd{ls} displays long format with permissions, owner, size, and modification time. The \texttt{-a} flag shows hidden files that begin with a dot. The \texttt{-h} flag presents sizes in human-readable units such as KB or MB.

Creating and managing files uses \cmd{mkdir} to create directories, \cmd{touch} to create empty files, \cmd{cp} to copy, \cmd{mv} to move or rename, and \cmd{rm} to delete. The \texttt{-p} flag for \cmd{mkdir} creates an entire directory path at once, including missing parent directories. The \texttt{-r} flag for \cmd{cp} and \cmd{rm} applies the operation recursively to directory contents.

To read file contents, \cmd{cat} displays all content at once, \cmd{head} displays the first few lines, and \cmd{tail} displays the last few lines. The \texttt{-n} flag for both commands sets the number of lines.

\begin{warningbox}
The \cmd{rm -rf} command deletes a directory and all its contents without confirmation. There is no trash folder. Verify the path before pressing Enter.
\end{warningbox}

\subsection*{Lab 2.3: Create a Practice Directory Structure}

The goal of this lab is to build a working directory structure and practice basic file manipulation.

\begin{enumerate}
    \item Move to this chapter's working directory:
\begin{lstlisting}[language=bash, caption={Entering the working directory}]
cd ~/lab-os/chapter02-hardware
pwd
\end{lstlisting}

    \item Create a nested directory structure in one command:
\begin{lstlisting}[language=bash, caption={Creating a directory structure}]
mkdir -p proyek/src proyek/docs proyek/logs
ls -l proyek/
\end{lstlisting}
    \textit{Observe:} The \texttt{-p} flag creates every directory in the path. Without \texttt{-p}, the command fails if the parent directory does not already exist.

    \item Create several example files and fill one of them:
\begin{lstlisting}[language=bash, caption={Creating and filling files}]
touch proyek/docs/readme.txt
echo "INFO: service started" > proyek/logs/app.log
echo "WARN: disk almost full" >> proyek/logs/app.log
echo "ERROR: connection failed" >> proyek/logs/app.log
cat proyek/logs/app.log
\end{lstlisting}

    \item Copy the log file to another directory:
\begin{lstlisting}[language=bash, caption={Copying a file}]
cp proyek/logs/app.log proyek/docs/app-backup.log
ls -lh proyek/docs/
\end{lstlisting}

    \item Rename the backup file:
\begin{lstlisting}[language=bash, caption={Renaming a file}]
mv proyek/docs/app-backup.log proyek/docs/app-2024.log
ls proyek/docs/
\end{lstlisting}

    \item Display the last three lines of the log:
\begin{lstlisting}[language=bash, caption={Reading the end of a log file}]
tail -3 proyek/logs/app.log
\end{lstlisting}
\end{enumerate}

\begin{challengebox}
Create the directory structure \texttt{archive/2024/01}, \texttt{archive/2024/02}, and \texttt{archive/2024/03} using \textbf{one} \cmd{mkdir -p} command with brace expansion. Check the result with \cmd{ls -R archive/}. Then delete the entire \texttt{archive} directory at once using \cmd{rm -r archive/}.
\end{challengebox}

%% ------------------------------------------------------------------------
\section{Text Manipulation}
\label{sec:text-manipulation}

Almost all data in Linux is text: system logs, configuration files, and command output. The ability to filter, modify, and extract information from text is a skill you will use every day as a system administrator.

The \cmd{grep} command searches for lines that match a pattern. The \texttt{-n} flag displays line numbers, \texttt{-i} ignores uppercase and lowercase differences, \texttt{-v} displays lines that do not match, and \texttt{-r} searches recursively through all files in a directory. To search for two patterns at once, use the \texttt{-E} flag with the pattern \texttt{pattern1|pattern2}.

The \cmd{sed} command is a stream editor that processes text line by line. Its most common operation is substitution with the syntax \texttt{s/old/new/g}: the letter \texttt{s} means substitution, and \texttt{g} at the end means replace every occurrence on a line. The \texttt{-i} flag applies changes directly to the original file. For safety on production systems, use \texttt{-i.bak} so the original file is saved as a backup.

The \cmd{awk} command processes text by columns. Each line is split into columns separated by spaces, and you can access a specific column with \texttt{\$1}, \texttt{\$2}, and so on. The variable \texttt{\$NF} refers to the last column. The variable \texttt{NR} contains the current line number. You can add a condition such as \texttt{NR>1} to skip a header line.

\subsection*{Lab 2.4: Filter and Extract Text from System Output}

The goal of this lab is to use \cmd{grep}, \cmd{sed}, and \cmd{awk} to process real system command output.

\begin{enumerate}
    \item Prepare a practice log file with more content:
\begin{lstlisting}[language=bash, caption={Creating a practice log file}]
cat > ~/lab-os/chapter02-hardware/sistem.log << 'EOF'
INFO: web service started on port 8080
WARN: CPU usage reached 75%
ERROR: failed to read configuration file
INFO: data backup completed
WARN: remaining memory below 20%
ERROR: database connection disconnected
INFO: system update scheduled
EOF
cat sistem.log
\end{lstlisting}

    \item Search for lines containing the word ERROR:
\begin{lstlisting}[language=bash, caption={Filtering ERROR lines}]
grep "ERROR" sistem.log
\end{lstlisting}
    \textit{Observe:} The result contains only lines with the word ERROR. Try adding the \texttt{-n} flag to show line numbers.

    \item Display lines that are not INFO:
\begin{lstlisting}[language=bash, caption={Displaying lines other than INFO}]
grep -v "INFO" sistem.log
\end{lstlisting}

    \item Replace \texttt{WARN} with \texttt{PERINGATAN} on all lines:
\begin{lstlisting}[language=bash, caption={Text substitution with sed}]
sed 's/WARN/PERINGATAN/g' sistem.log
\end{lstlisting}
    \textit{Notice:} The original file does not change. Only the output shown in the terminal has been modified.

    \item Create a practice configuration file and change its value:
\begin{lstlisting}[language=bash, caption={Editing configuration with sed}]
cat > server.conf << 'EOF'
PORT=8080
MODE=development
HOST=localhost
EOF
sed -i 's/development/production/g' server.conf
cat server.conf
\end{lstlisting}

    \item Extract the first and last columns from \cmd{df -h} output:
\begin{lstlisting}[language=bash, caption={Extracting columns with awk}]
df -h | awk '{print $1, $NF}'
\end{lstlisting}
    \textit{Observe:} \texttt{\$1} is the first column, the filesystem name, and \texttt{\$NF} is the last column, the mount point. The output is much shorter than the full \cmd{df -h} output.

    \item Filter filesystems with disk usage above 50\%:
\begin{lstlisting}[language=bash, caption={Filtering with an awk condition}]
df -h | awk 'NR==1 {print} NR>1 && ($5+0) > 50 {print}'
\end{lstlisting}
\end{enumerate}

\begin{challengebox}
Combine \cmd{grep} and \cmd{awk} in one pipeline to extract information from \file{/proc/meminfo}. Find lines containing \texttt{Mem} with \cmd{grep}, then extract the field name, the first column, and its value, the second column, with \cmd{awk}. The command is: \texttt{grep Mem /proc/meminfo | awk '\{print \$1, \$2, \$3\}'}. Add \cmd{sort} at the end of the pipeline to sort the result.
\end{challengebox}

%% ------------------------------------------------------------------------
\section{Command Help and Documentation}
\label{sec:command-documentation}

Linux has many commands, and no administrator memorizes all of them. The more important skill is knowing how to find answers quickly. Linux provides rich built-in documentation: manual pages, quick option summaries through \cmd{--help}, GNU documentation through \cmd{info}, and topic indexes through \cmd{apropos}.

The \cmd{man command-name} command displays a complete manual page containing descriptions, options, examples, and related files. Almost all modern utilities also support \cmd{command-name --help} for a quick summary. For some GNU programs, \cmd{info command-name} provides longer and more structured documentation. Meanwhile, \cmd{apropos keyword} is useful when you do not yet know the command name, for example when searching for commands related to ``copy'' or ``archive''.

There are also other helper commands often used in the terminal: \cmd{which} to find an executable location in \var{PATH}, \cmd{type} to determine whether a name is a shell builtin, alias, function, or regular executable, and \cmd{help} for help on Bash built-in commands.

This approach matters for this book: instead of stopping when you meet a new command, get used to opening its documentation, reading the relevant options, and trying a small example directly in the terminal.

\subsection*{Lab 2.5: Explore Command Documentation}

The goal of this lab is to become comfortable finding documentation and understanding the different help sources available in Linux.

\begin{enumerate}
    \item Open the manual page for \cmd{ls} and search for the section explaining the \texttt{-a} flag:
\begin{lstlisting}[language=bash, caption={Opening a manual page}]
man ls
\end{lstlisting}
    \textit{Observe:} Inside \cmd{man}, use the \texttt{/} key to search for a keyword such as \texttt{-a}, then press \texttt{n} for the next result and \texttt{q} to quit.

    \item View quick option summaries for \cmd{grep} and \cmd{tar}:
\begin{lstlisting}[language=bash, caption={Quick help summaries}]
grep --help | head -20
tar --help | head -20
\end{lstlisting}

    \item Use \cmd{apropos} to search for commands by topic:
\begin{lstlisting}[language=bash, caption={Searching for commands by keyword}]
apropos copy
apropos archive
\end{lstlisting}

    \item Check the location and type of the \cmd{cd}, \cmd{ls}, and \cmd{grep} commands:
\begin{lstlisting}[language=bash, caption={Identifying command types in the shell}]
type cd
type ls
which grep
\end{lstlisting}

    \item Read help for a shell builtin and info documentation if available:
\begin{lstlisting}[language=bash, caption={Builtin help and info}]
help cd
info coreutils 'ls invocation'
\end{lstlisting}
\end{enumerate}

\begin{challengebox}
Choose three commands from this chapter, for example \cmd{ls}, \cmd{grep}, and \cmd{tar}. For each one, record one way to open full documentation, one way to view brief help, and one option you find most useful. Explain when you would use that option.
\end{challengebox}

%% ------------------------------------------------------------------------
\section{Summary}

In this chapter, you learned:

\begin{itemize}[noitemsep,topsep=2pt]
    \item \textbf{Hardware detection}: \cmd{lspci}, \cmd{lsusb}, \cmd{lscpu}, and \cmd{dmesg} show system components detected by the kernel.
    \item \textbf{Kernel modules}: \cmd{lsmod} displays active modules, \cmd{modinfo} reads module details, and \cmd{modprobe} loads a module with its dependencies.
    \item \textbf{Basic terminal commands}: \cmd{ls}, \cmd{cd}, \cmd{pwd}, \cmd{mkdir -p}, \cmd{cp}, \cmd{mv}, \cmd{rm}, \cmd{cat}, \cmd{head}, and \cmd{tail} are the foundation of navigation and file manipulation.
    \item \textbf{Text manipulation}: \cmd{grep} filters lines by pattern, \cmd{sed} replaces text, and \cmd{awk} extracts columns from formatted output.
    \item \textbf{Command documentation}: \cmd{man}, \cmd{--help}, \cmd{info}, \cmd{apropos}, \cmd{type}, and \cmd{which} help you learn new commands independently.
\end{itemize}

%% ------------------------------------------------------------------------
\section{Lab Assignment}

\textbf{General Instructions}: Complete all exercises independently. Record important steps, save screenshots when needed, and summarize the results as personal documentation.

\subsubsection*{Lab Assignment 2.1 VM Hardware Profile}
Run \cmd{lspci -nnk} and choose one PCI device. Record: (a) the device name, (b) the vendor and device ID in hexadecimal format, (c) the active kernel \term{driver} name. Then run \cmd{modinfo} for that \term{driver} and write a brief description and the module author's name.


\subsubsection*{Lab Assignment 2.2 Documentation and Kernel Modules}
Choose one active module from the \cmd{lsmod} output. Document: (a) the module name, (b) the description and license from \cmd{modinfo}, (c) one related kernel message from \cmd{dmesg}, if any. Then use \cmd{man} or \cmd{--help} to find the meaning of one option from the \cmd{modprobe} command and explain its function.


\subsubsection*{Lab Assignment 2.3 Log Analysis with grep/sed/awk}
Create a \file{server.log} file containing at least 10 lines with mixed message levels: INFO, WARN, and ERROR. Then: (a) use \cmd{grep -c ERROR server.log} to count ERROR lines, (b) use \cmd{sed} to replace every \texttt{WARN} with \texttt{WARNING} and save it to a new file named \file{server-clean.log}, (c) use \cmd{awk} to display only the first word, the message level, from each line. Include screenshots of each command output.


%
